\documentclass[10pt,a4paper]{amsart}
\usepackage[margin=19mm]{geometry}
\usepackage{amsmath,amssymb,mathtools}
\usepackage[scr=rsfs,cal=euler]{mathalfa}
\usepackage{xfrac}
\usepackage[unicode,hidelinks,bookmarks=false]{hyperref}
\newcommand{\ie}{i.e.\ }
\newcommand{\cf}{cf.\ }
\newcommand{\resp}{respectively\ }
\newcommand{\Subsection}{\subsection}
\providecommand{\nobreakdash}{\nobreak\hbox{-}}
\setlength{\emergencystretch}{2em}
\multlinegap0pt
\usepackage{bm}
\usepackage{bbm}
\usepackage{dsfont}
\usepackage{xspace}
\usepackage{cancel}
\usepackage{mathtools}
\usepackage{amsthm}
\makeatletter
\@ifundefined{theorem}{\newtheorem{theorem}{Theorem}}{}
\makeatother
\title[On characteristic foliations of metric contact-symplectic st]{On characteristic foliations of metric contact-symplectic structures}
\author{Amine Hadjar}
\date{Source: arXiv:2602.03356v1 (CC BY 4.0)}
\begin{document}
\maketitle

\noindent\emph{Source and licence.} On characteristic foliations of metric contact-symplectic structures. Version arXiv:2602.03356v1 (CC BY 4.0), distributed under the Creative Commons Attribution 4.0 International licence, as recorded in the arXiv licence metadata for this exact version and shown on the abstract page https://arxiv.org/abs/2602.03356v1. Full rights notice: licenses/arxiv-2602.03356-CC-BY-4.0.txt.\par\smallskip

\noindent\textit{Editorial scope.} Counted unit: the Theorem whose source label is \texttt{theoremdV} in the article (source0.tex lines 352--357, source label \texttt{theoremdV}), delivered together with the article's own complete original proof of it (source0.tex lines 358--386, one \texttt{proof} environment of 29 lines). No mathematics is written, repaired or re-derived: statement and proof are the article's own text line for line. Required context retained uncounted: none. Every definition, display and label of those lines is retained unchanged. Omitted from this package and not redistributed: 1--351, 387--475. Nothing inside a retained excerpt is omitted or altered: every retained body is delivered exactly as the article's lines are written, so no omission receipt of any kind accompanies this package. Citations: the retained excerpts carry no citation command at all, so no bibliography entry of the article is transcribed and no reference is redistributed. Reference adaptation: none was needed and none was made; every reference inside the delivered unit resolves inside it, so no unresolved reference or citation is produced. Printed slips kept literal: no printed word, symbol, hypothesis or number of the counted statement or of its proof was altered. Layout and class: the article's own class is \texttt{amsart} with packages including amscd, amssymb, times, amsmath, color, fullpage, enumitem; the wrapper is the lane's pinned \texttt{amsart} editorial wrapper at 10pt on A4, which restates the theorem-like environments the retained text needs and adds the article's own macro definitions that the retained text needs (none), with extra wrapper packages (bm, bbm, dsfont, xspace, cancel, mathtools). The delivered numbering is the unit's own. Assets: no retained span contains a figure environment, a graphics-inclusion command, a PDF-page inclusion, a table or a TikZ picture; no image file is copied into this package, the package declares no asset dependency and no image file is redistributed with it. Licence: version arXiv:2602.03356v1 of this article is distributed under the Creative Commons Attribution 4.0 International licence, as recorded in the arXiv licence metadata for this exact version and shown on the abstract page https://arxiv.org/abs/2602.03356v1. A full rights notice is delivered at licenses/arxiv-2602.03356-CC-BY-4.0.txt. Characters: every retained excerpt is pure ASCII, so the delivered engine, XeLaTeX with its default Latin Modern text font, reports no missing character.
\begin{theorem}\label{theoremdV}
For a contact-symplectic pair $(\eta, \omega)$ of type $(m,n)$ on a manifold, all associated metrics have the same volume element given by
\begin{equation*}
dV= \dfrac{(-1)^{m+n}}{2^{m+n} m! \ n!}\ \eta\wedge (d\eta)^{m}\wedge\omega^{n}
\end{equation*}
\end{theorem}

\begin{proof}
We know that for a metric contact-symplectic structure $(\eta, \omega, g, \phi)$ of type $(m,n)$ with Reeb vector field $\xi$ on a manifold $M$, $(\phi, \xi, \eta, g)$ is an almost contact metric structure on $M$.
Consequently, one can locally construct a $\phi$-basis, that is, an orthonormal frame $\{ \xi, X_1, \dots , X_{m+n},  X_{1^*},$ $\dots , X_{(m+n)^*}\}$ where $X_{i^*}=\phi X_i$.
Let $\{ \eta, \theta^1, \dots , \theta^{m+n}, \theta^{1^*}, \dots , \theta^{(m+n)^*}\}$ be its dual basis. 
Then
\begin{equation*}
dV= \eta \wedge \theta^1 \wedge  \theta^{1^*} \wedge  \dots \wedge \theta^{m+n} \wedge \theta^{(m+n)^*}.
\end{equation*}
Observe that
\begin{equation*}
d\eta+\omega =\sum_{i=1}^{m+n} (\theta^{i^*} \wedge \theta^i-\theta^i \wedge  \theta^{i^*})=-2\sum_{i=1}^{m+n}\theta^i \wedge  \theta^{i^*}
\end{equation*}
which follows from the identity $g(X, \phi Y) = (\mathrm{d}\eta + \omega)(X, Y)$ and the orthonormality of the local $\phi$-basis. 
Next we have
\begin{equation*}
(d\eta+\omega )^{m+n}=(-2)^{m+n} (m+n)! \ \theta^1 \wedge  \theta^{1^*} \wedge  \dots \wedge \theta^{m+n} \wedge \theta^{(m+n)^*}.
\end{equation*}
On the other hand,
\begin{align*}
(d\eta+\omega )^{m+n}&=\sum_{k=0}^{m+n} \binom{m+n}{k} (d\eta)^k \wedge \omega^{m+n-k} \\
				   &= \dfrac{(m+n)!}{m!\ n!}(d\eta)^m \wedge \omega^n .
\end{align*}
Comparing both expressions yields:
\begin{align*}
\eta\wedge(d\eta)^m \wedge \omega^n&=(-2)^{m+n} m!\ n! \ \eta \wedge \theta^1 \wedge  \theta^{1^*} \wedge  \dots \wedge \theta^{m+n} \wedge \theta^{(m+n)^*}\\
								&=(-2)^{m+n} m!\ n!  \ dV .
\end{align*}
from which the formula follows.
\end{proof}

\end{document}
